% ============================================================
%  PREAMBLE — Beamer-Präsentation
%  Thema: Komplexe Zahlen
% ============================================================

% ---- Dokumentklasse --------------------------------------------
\documentclass[
  aspectratio=169,
  10pt
]{beamer}

% ---- Sprache ---------------------------------------------------
\usepackage[ngerman]{babel}
\usepackage[utf8]{inputenc}
\usepackage[T1]{fontenc}

% Zeichen aus dem Arabischen Transliterations-Standard (für Bib-Einträge)
\DeclareUnicodeCharacter{0101}{\={a}}
\DeclareUnicodeCharacter{0100}{\={A}}
\DeclareUnicodeCharacter{012B}{\={i}}
\DeclareUnicodeCharacter{012A}{\={I}}
\DeclareUnicodeCharacter{016B}{\={u}}
\DeclareUnicodeCharacter{016A}{\={U}}
\DeclareUnicodeCharacter{1E25}{\d{h}}
\DeclareUnicodeCharacter{1E24}{\d{H}}
\DeclareUnicodeCharacter{1E63}{\d{s}}
\DeclareUnicodeCharacter{1E62}{\d{S}}
\DeclareUnicodeCharacter{00E0}{\`{a}}

% ---- Beamer-Theme ----------------------------------------------
\usetheme{Madrid}
\usecolortheme{default}
\usefonttheme{professionalfonts}
\setbeamertemplate{navigation symbols}{}
\setbeamertemplate{theorems}[numbered]

% Farben anpassen
\definecolor{mainblue}{RGB}{0,70,127}
\definecolor{accentred}{RGB}{180,20,20}
\setbeamercolor{palette primary}{bg=mainblue, fg=white}
\setbeamercolor{palette secondary}{bg=mainblue!80, fg=white}
\setbeamercolor{palette tertiary}{bg=mainblue!60, fg=white}
\setbeamercolor{palette quaternary}{bg=mainblue, fg=white}
\setbeamercolor{frametitle}{bg=mainblue, fg=white}
\setbeamercolor{title}{fg=white}
\setbeamercolor{structure}{fg=mainblue}
\setbeamercolor{block title}{bg=mainblue!90, fg=white}
\setbeamercolor{block body}{bg=mainblue!10}
\setbeamercolor{block title example}{bg=accentred!80, fg=white}
\setbeamercolor{block body example}{bg=accentred!10}
\setbeamercolor{alerted text}{fg=accentred}

% ---- Mathematik-Pakete -----------------------------------------
\usepackage{amsmath, amsfonts, amssymb, amsthm}
\usepackage{mathtools}
\usepackage{bbm}

\newcommand\scalemath[2]{\scalebox{#1}{\mbox{\ensuremath{\displaystyle #2}}}}


% ---- TikZ & PGF ------------------------------------------------
\usepackage{tikz}
\usepackage{tikz-cd}
\usepackage{pgfplots}
\pgfplotsset{compat=newest}
\usetikzlibrary{arrows.meta, calc, angles, quotes, backgrounds, patterns, decorations.pathmorphing, babel}
\tikzset{>=latex}

% ---- Weitere Pakete --------------------------------------------
\usepackage{graphicx}
\usepackage{booktabs}
\usepackage{csquotes}
\usepackage{xcolor}

% ---- Bibliographie ---------------------------------------------
\usepackage[backend=biber, style=alphabetic]{biblatex}
\addbibresource{literatur.bib}

\DeclareCiteCommand{\cite}[\mkbibbrackets]
  {\usebibmacro{prenote}}
  {\printnames{labelname},\addspace
   \printfield{title},\addspace
   \printfield{year}}
  {\multicitedelim}
  {\usebibmacro{postnote}}

% ============================================================
%  THEOREM-UMGEBUNGEN (Beamer-Stil)
%  Beamer definiert bereits: theorem, lemma, corollary,
%  definition, example, examples, fact, proof.
%  Hier nur deutsche Aliases und neue Umgebungen.
% ============================================================

\newtheorem{satz}[theorem]{Satz}
\newtheorem{korollar}[theorem]{Korollar}
\newtheorem{beispiel}[theorem]{Beispiel}
\newtheorem{bemerkung}[theorem]{Bemerkung}

% ============================================================
%  PRETTY SYMBOLS
% ============================================================
\let\uglyepsilon\epsilon  \let\epsilon\varepsilon
\let\uglyphi\phi          \let\phi\varphi


% ============================================================
%  MAKROS
% ============================================================

% ---- Zahlkörper -----------------------------------------------
\providecommand{\R}{\mathbb{R}}
\providecommand{\N}{\mathbb{N}}
\providecommand{\Z}{\mathbb{Z}}
\providecommand{\C}{\mathbb{C}}
\providecommand{\Q}{\mathbb{Q}}

% ---- Allgemeine Operatoren ------------------------------------
\DeclarePairedDelimiter\abs{\lvert}{\rvert}
\providecommand{\norm}[1]{\left\lVert#1\right\rVert}
\providecommand{\ii}{\mathrm{i}}
\providecommand{\euler}{\mathrm{e}}
\providecommand{\diff}{\mathop{}\!\mathrm{d}}
\renewcommand{\Re}{\operatorname{Re}}
\renewcommand{\Im}{\operatorname{Im}}
\providecommand{\Arg}{\operatorname{Arg}}
\providecommand{\conj}[1]{\overline{#1}}
\newcommand{\mapsfrom}{\mathrel{\reflectbox{\ensuremath{\mapsto}}}}

% ---- Hilfskommandos -------------------------------------------
\newcommand{\highlight}[1]{\textcolor{accentred}{\textbf{#1}}}
\newcommand{\emphmath}[1]{\textcolor{mainblue}{#1}}
% Polarkoordinaten-Tupel: farbig unterstrichen zur klaren Unterscheidung
\newcommand{\polar}[1]{\textcolor{mainblue}{\underline{\left(#1\right)}}}
